\section{Sémantická analýza}

Sémantická analýza nie je samostatný prechod. Prebieha priebežne \textbf{počas syntaktickej analýzy}: parser pri spracovaní každého príkazu a pri každej redukcii vo výraze volá funkcie modulov \texttt{symtable}, \texttt{semantics} a \texttt{builtins}, takže už vytvorený AST je typovo anotovaný (každý uzol má vyplnené pole \texttt{type}).

\subsection{Tabuľka symbolov}
Tabuľka symbolov (\texttt{symtable.c}) je hašovacia tabuľka s reťazením synonym (funkcia DJB2, 101 záznamov). Každý rozsah platnosti má vlastnú tabuľku s odkazom na rodičovský rozsah, vyhľadávanie (\texttt{symtable\_lookup}) postupuje od aktuálneho rozsahu ku globálnemu, kým \texttt{symtable\_lookup\_local} hľadá iba v aktuálnom. Záznam obsahuje názov, druh (\texttt{SYMBOL\_VAR}, \texttt{SYMBOL\_VAL}, \texttt{SYMBOL\_FUNC}), dátový typ, príznak zmeniteľnosti a pre funkcie počet a typy parametrov. Vloženie už existujúceho názvu v tom istom rozsahu vráti \texttt{ERR\_SEMANTIC\_REDEFINITION}.

Funkcia sa do globálnej tabuľky vloží ešte pred spracovaním parametrov a tela. Jej signatúra (návratový typ, počet a typy parametrov) sa doplní hneď po prečítaní hlavičky, takže rekurzívne volania sa kontrolujú rovnako ako iné. Parametre sú uložené ako nezmeniteľné premenné.

\subsection{Vstavané funkcie}
Modul \texttt{builtins.c} pred spustením parsera zaregistruje do globálneho rozsahu funkcie jazyka IFJ26 s ich signatúrami: \texttt{readLine}, \texttt{readInt}, \texttt{readDouble}, \texttt{print}, \texttt{toStr}, \texttt{toInt}, \texttt{strLength}, \texttt{ord}, \texttt{chr} a \texttt{subStr}. Parameter typu \texttt{TYPE\_UNKNOWN} (pri \texttt{print} a \texttt{toStr}) zastupuje ``ľubovoľný typ''. Pamäť s typmi parametrov vlastní záznam a uvoľní sa v \texttt{symtable\_free}.

\subsection{Typový systém}
\begin{itemize}
    \item \textbf{Kompatibilita priradenia} (\texttt{is\_type\_compatible}): rovnaké typy, alebo \texttt{Int} do \texttt{Double}. Neznámy typ (\texttt{TYPE\_UNKNOWN}) je kompatibilný so všetkým, aby sa jedna chyba nešírila ďalej.
    \item \textbf{Aritmetika:} \texttt{+ - * /} nad \texttt{Int}/\texttt{Double}; ak je aspoň jeden operand \texttt{Double}, výsledok je \texttt{Double}. Operátor \texttt{+} je navyše definovaný nad dvojicou \texttt{String}.
    \item \textbf{Relačné operátory} \texttt{< <= > >=} sú definované nad číslami, výsledok je \texttt{Boolean}.
    \item \textbf{Rovnosť} \texttt{== !=} vyžaduje rovnaké (alebo obe číselné) typy, výsledok je \texttt{Boolean}.
    \item \textbf{Logické operátory} \texttt{\&\&}, \texttt{||} a \texttt{!} vyžadujú \texttt{Boolean}. Unárne \texttt{-} vyžaduje číslo.
    \item \textbf{Podmienky} v \texttt{if}, \texttt{while} a stráže cyklu \texttt{for} musia byť typu \texttt{Boolean}; hranice a krok cyklu \texttt{for} typu \texttt{Int}.
    \item \textbf{\texttt{if} ako výraz:} ak sa typy oboch vetiev líšia, ale obe sú číselné, výsledkom je \texttt{Double}; inak musia byť zhodné.
    \item \textbf{Odvodenie typu:} pri \texttt{val}/\texttt{var} bez uvedeného typu sa typ premennej odvodí z inicializačného výrazu.
\end{itemize}

\subsection{Kontroly a chybové kódy}
\begin{longtable}{p{0.8cm} p{6.3cm} p{7.5cm}}
\toprule
\textbf{Kód} & \textbf{Konštanta} & \textbf{Kedy nastane} \\
\midrule
\endhead
3 & \texttt{ERR\_SEMANTIC\_UNDEFINED} & použitie nedefinovanej premennej alebo funkcie; chýbajúca funkcia s anotáciou \texttt{@main} \\
4 & \texttt{ERR\_SEMANTIC\_REDEFINITION} & opätovná definícia funkcie alebo premennej v tom istom rozsahu \\
5 & \texttt{ERR\_SEMANTIC\_FUNCTION\_CALL} & zlý počet alebo typ argumentov volania; parametre pri funkcii \texttt{@main} \\
6 & \texttt{ERR\_SEMANTIC\_COMPATIBILITY} & nekompatibilný typ pri inicializácii, priradení, \texttt{return} alebo vetvách \texttt{if} výrazu \\
7 & \texttt{ERR\_SEMANTIC\_FORBIDDEN\_WRITING} & zápis do \texttt{val}, parametra funkcie alebo funkcie \\
8 & \texttt{ERR\_SEMANTIC\_TYPE} & nepovolené typy operandov operátora alebo nebooleovská podmienka \\
\bottomrule
\end{longtable}

Funkcia \texttt{validate\_function\_call} overí existenciu symbolu, že ide o funkciu, počet argumentov a kompatibilitu typu každého argumentu, a vráti návratový typ funkcie. Rovnakú kontrolu používa volanie cez bodku (\texttt{x.f(a)}), kde sa prijímač počíta ako prvý argument. Funkcia \texttt{check\_variable\_assignment} overí, že cieľ priradenia je zmeniteľná premenná (\texttt{var}). Pri návrate z funkcie (\texttt{return} aj jednovýrazové telo) sa typ výrazu porovnáva s deklarovaným návratovým typom.

\subsection{Známe obmedzenia}
Spracovanie je jednoprechodové, preto funkciu možno volať až po jej definícii v zdrojovom texte. Jazyk zatiaľ nevyužíva výsledky tejto analýzy na generovanie kódu (modul \texttt{codegen} nie je implementovaný).